\subsection*{A.1 Prompts}
\textbf{Decision APIs.} Instruction (identical for N and D, all datasets): \emph{``Classify the user's message into the single most
appropriate intent.''} Options are the class identifiers with a description: in condition N the identifier with underscores replaced
by spaces, in condition D the description in \nolinkurl{data/descriptions/*_D.json} (Ottavo: \nolinkurl{data/ottavo/labels.json}).
In-context examples are appended as lines \texttt{- "<text>" -> <label>}. OpenAI receives the options as a \texttt{choice} question;
Baseten System One receives them as \texttt{criteria}. Full request bodies are reproducible from \nolinkurl{dmstudy/decisions.py}.

\textbf{Generalist controls.} One user turn with the instruction \emph{``Classify the customer's message into exactly one of the categories
below, based on the customer's main intent. Answer with the letter of the category only.''}, the categories listed as
\texttt{A) name: description}, and the message; chat template with thinking disabled. The next-token distribution is restricted to
the valid letters and renormalised; the mean probability mass on valid letters was close to 1.

\subsection*{A.2 Supervised configurations}
TF-IDF: word 1--2-grams and character 2--5-grams, logistic regression, $C\in\{1,10,100\}$ by 3-fold CV inside the budget (validation
split for full data). SetFit: \nolinkurl{intfloat/multilingual-e5-base}, prefix ``query: '', 256 tokens, 20 pair iterations, one epoch,
batch 16, logistic head, at most 2{,}000 contrastive steps for full data. mmBERT-base: full fine-tuning, learning rate $\in\{2\cdot10^{-5},
5\cdot10^{-5}\}$. Qwen3 0.6B/1.7B: LoRA $r=16$, $\alpha=32$, dropout 0.05 on all attention and MLP projections, linear classification head
trained in fp32, learning rate $\in\{5\cdot10^{-5}, 2\cdot10^{-4}\}$, batch 16. Qwen3-8B: QLoRA (NF4, double quantisation), same adapters,
learning rate $10^{-4}$, batch 8. Durations evaluated in CV: 5, 10, 20 epochs at 5 per class; 3, 6, 12 at 20 per class; full data: best of
4 (public) or 6 (Ottavo) epochs on the validation split. Maximum input length 256 tokens. Selected recipes are in \nolinkurl{results/selection/}.

\subsection*{A.3 Data}
Few-shot subsets, validation splits, API subsets and option permutations are listed by example ID in \nolinkurl{data/manifests/}. The
Ottavo generation prompts and the validators' instructions are in \nolinkurl{data/ottavo/generation/}; validator agreement and dropped
items are in \nolinkurl{reports/audit_initial.md} and \nolinkurl{data/ottavo_r/build_report.json}.

\subsection*{A.4 Statistics}
Paired percentile bootstrap over test messages, 2{,}000 replicates, seed 20261024, implemented with multinomial weights; macro-F1 over
classes present in the gold labels; supervised statistics are means over seeds within each replicate. Holm's correction is applied
within each family of contrasts (P1--P6 and S1 declared in advance, P2N added post hoc) listed in \nolinkurl{reports/contrasts_primary.md}. ECE uses 15 equal-width bins on the maximum class probability;
Brier is $\frac{1}{n}\sum_i\sum_k (p_{ik}-y_{ik})^2$. Selective prediction thresholds are chosen on one random half of the test set and
evaluated on the other (two-fold cross-fitting).

\subsection*{A.5 Per-class results}
\label{app:perclass}
Table~\ref{tab:perclass} summarises per-class F1 (supported classes only; compact classifiers: mean over three seeds). The complete
per-class table with supports is in \nolinkurl{reports/per_class_f1.csv}.

\subsection*{A.6 Variability across seeds}
Table~\ref{tab:seeds} reports, for each supervised configuration, the mean, standard deviation, minimum and maximum macro-F1 over the
three seeds (17, 42, 123). The seed changes both the few-shot subset and the training randomness. This run-to-run variability is distinct
from the uncertainty due to the finite test sample, which the bootstrap intervals in the main text describe.

\begin{table}[!htbp]
\centering\scriptsize
\caption{Macro-F1 across the three seeds.}
\label{tab:seeds}
\begin{tabular}{lllrrrr}
\toprule
Data & System & Examples/class & Mean & SD & Min & Max \\
\midrule
Banking77 & TF-IDF + LR & 5 & 0.645 & 0.016 & 0.626 & 0.655 \\
Banking77 & TF-IDF + LR & 20 & 0.820 & 0.003 & 0.818 & 0.823 \\
Banking77 & TF-IDF + LR & all & 0.900 & 0.001 & 0.900 & 0.901 \\
Banking77 & SetFit (mE5-base) & 5 & 0.717 & 0.010 & 0.709 & 0.729 \\
Banking77 & SetFit (mE5-base) & 20 & 0.839 & 0.017 & 0.825 & 0.857 \\
Banking77 & SetFit (mE5-base) & all & 0.864 & 0.003 & 0.860 & 0.866 \\
Banking77 & mmBERT-base & 5 & 0.547 & 0.015 & 0.537 & 0.565 \\
Banking77 & mmBERT-base & 20 & 0.835 & 0.010 & 0.824 & 0.843 \\
Banking77 & mmBERT-base & all & 0.904 & 0.004 & 0.900 & 0.908 \\
Banking77 & Qwen3-0.6B LoRA & 5 & 0.575 & 0.022 & 0.551 & 0.594 \\
Banking77 & Qwen3-0.6B LoRA & 20 & 0.846 & 0.033 & 0.808 & 0.871 \\
Banking77 & Qwen3-0.6B LoRA & all & 0.910 & 0.004 & 0.907 & 0.915 \\
\midrule
MASSIVE-it & TF-IDF + LR & 5 & 0.491 & 0.013 & 0.482 & 0.507 \\
MASSIVE-it & TF-IDF + LR & 20 & 0.666 & 0.003 & 0.664 & 0.669 \\
MASSIVE-it & TF-IDF + LR & all & 0.786 & 0.005 & 0.781 & 0.789 \\
MASSIVE-it & SetFit (mE5-base) & 5 & 0.615 & 0.015 & 0.599 & 0.628 \\
MASSIVE-it & SetFit (mE5-base) & 20 & 0.688 & 0.017 & 0.670 & 0.704 \\
MASSIVE-it & SetFit (mE5-base) & all & 0.696 & 0.017 & 0.684 & 0.716 \\
MASSIVE-it & mmBERT-base & 5 & 0.447 & 0.019 & 0.426 & 0.462 \\
MASSIVE-it & mmBERT-base & 20 & 0.663 & 0.004 & 0.661 & 0.668 \\
MASSIVE-it & mmBERT-base & all & 0.803 & 0.004 & 0.800 & 0.807 \\
MASSIVE-it & Qwen3-0.6B LoRA & 5 & 0.358 & 0.035 & 0.318 & 0.380 \\
MASSIVE-it & Qwen3-0.6B LoRA & 20 & 0.667 & 0.008 & 0.658 & 0.673 \\
MASSIVE-it & Qwen3-0.6B LoRA & all & 0.795 & 0.009 & 0.785 & 0.801 \\
\midrule
Ottavo & TF-IDF + LR & 5 & 0.588 & 0.011 & 0.576 & 0.598 \\
Ottavo & TF-IDF + LR & 20 & 0.814 & 0.006 & 0.807 & 0.818 \\
Ottavo & TF-IDF + LR & all & 0.915 & 0.000 & 0.915 & 0.915 \\
Ottavo & SetFit (mE5-base) & 5 & 0.766 & 0.013 & 0.757 & 0.781 \\
Ottavo & SetFit (mE5-base) & 20 & 0.907 & 0.013 & 0.893 & 0.919 \\
Ottavo & SetFit (mE5-base) & all & 0.950 & 0.004 & 0.948 & 0.955 \\
Ottavo & mmBERT-base & 5 & 0.502 & 0.034 & 0.475 & 0.541 \\
Ottavo & mmBERT-base & 20 & 0.833 & 0.022 & 0.810 & 0.854 \\
Ottavo & mmBERT-base & all & 0.916 & 0.008 & 0.909 & 0.924 \\
Ottavo & Qwen3-0.6B LoRA & 5 & 0.150 & 0.026 & 0.127 & 0.178 \\
Ottavo & Qwen3-0.6B LoRA & 20 & 0.764 & 0.041 & 0.724 & 0.807 \\
Ottavo & Qwen3-0.6B LoRA & all & 0.913 & 0.009 & 0.906 & 0.923 \\
Ottavo & Qwen3-1.7B LoRA & 5 & 0.236 & 0.047 & 0.201 & 0.289 \\
Ottavo & Qwen3-1.7B LoRA & 20 & 0.820 & 0.028 & 0.801 & 0.852 \\
Ottavo & Qwen3-1.7B LoRA & all & 0.950 & 0.004 & 0.948 & 0.955 \\
Ottavo & Qwen3-8B QLoRA & 5 & 0.277 & 0.002 & 0.275 & 0.280 \\
Ottavo & Qwen3-8B QLoRA & 20 & 0.893 & 0.020 & 0.881 & 0.917 \\
Ottavo & Qwen3-8B QLoRA & all & 0.975 & 0.002 & 0.973 & 0.977 \\
\bottomrule
\end{tabular}
\end{table}

\begin{table}[!htbp]
\centering\footnotesize
\caption{Per-class F1: median, 10th percentile, minimum, number of classes below 0.5 and the three weakest classes (all data for
compact classifiers; zero-shot with descriptions for decision models).}
\label{tab:perclass}
\begin{tabular}{llrrrr>{\raggedright\arraybackslash}p{4.9cm}}
\toprule
Data & System & Median & 10th pct. & Min & $<0.5$ & Three weakest classes \\
\midrule
Banking77 & Mammoth-1 & 0.929 & 0.815 & 0.692 & 0/77 & topping\_\allowbreak{}up\_\allowbreak{}by\_\allowbreak{}card, transfer\_\allowbreak{}not\_\allowbreak{}received\_\allowbreak{}by\_\allowbreak{}recipient, exchange\_\allowbreak{}via\_\allowbreak{}app \\
Banking77 & Elephant-1 & 0.929 & 0.818 & 0.690 & 0/77 & topping\_\allowbreak{}up\_\allowbreak{}by\_\allowbreak{}card, card\_\allowbreak{}arrival, declined\_\allowbreak{}transfer \\
Banking77 & OpenAI Decisions & 0.857 & 0.712 & 0.353 & 3/77 & top\_\allowbreak{}up\_\allowbreak{}by\_\allowbreak{}bank\_\allowbreak{}transfer\_\allowbreak{}charge, topping\_\allowbreak{}up\_\allowbreak{}by\_\allowbreak{}card, beneficiary\_\allowbreak{}not\_\allowbreak{}allowed \\
Banking77 & TF-IDF + LR & 0.917 & 0.800 & 0.750 & 0/77 & top\_\allowbreak{}up\_\allowbreak{}failed, balance\_\allowbreak{}not\_\allowbreak{}updated\_\allowbreak{}after\_\allowbreak{}bank\_\allowbreak{}transfer, reverted\_\allowbreak{}card\_\allowbreak{}payment? \\
Banking77 & SetFit (mE5-base) & 0.892 & 0.746 & 0.450 & 2/77 & get\_\allowbreak{}disposable\_\allowbreak{}virtual\_\allowbreak{}card, virtual\_\allowbreak{}card\_\allowbreak{}not\_\allowbreak{}working, disposable\_\allowbreak{}card\_\allowbreak{}limits \\
Banking77 & mmBERT-base & 0.917 & 0.809 & 0.744 & 0/77 & lost\_\allowbreak{}or\_\allowbreak{}stolen\_\allowbreak{}card, top\_\allowbreak{}up\_\allowbreak{}by\_\allowbreak{}bank\_\allowbreak{}transfer\_\allowbreak{}charge, balance\_\allowbreak{}not\_\allowbreak{}updated\_\allowbreak{}after\_\allowbreak{}bank\_\allowbreak{}transfer \\
Banking77 & Qwen3-0.6B LoRA & 0.913 & 0.826 & 0.788 & 0/77 & topping\_\allowbreak{}up\_\allowbreak{}by\_\allowbreak{}card, declined\_\allowbreak{}card\_\allowbreak{}payment, declined\_\allowbreak{}transfer \\
\midrule
MASSIVE-it & Mammoth-1 & 0.900 & 0.667 & 0.000 & 2/59 & iot\_\allowbreak{}hue\_\allowbreak{}lighton, music\_\allowbreak{}settings, general\_\allowbreak{}greet \\
MASSIVE-it & Elephant-1 & 0.905 & 0.667 & 0.500 & 0/59 & datetime\_\allowbreak{}convert, iot\_\allowbreak{}wemo\_\allowbreak{}on, music\_\allowbreak{}settings \\
MASSIVE-it & OpenAI Decisions & 0.827 & 0.557 & 0.000 & 5/59 & iot\_\allowbreak{}hue\_\allowbreak{}lighton, music\_\allowbreak{}dislikeness, general\_\allowbreak{}quirky \\
MASSIVE-it & TF-IDF + LR & 0.866 & 0.511 & 0.000 & 3/59 & iot\_\allowbreak{}hue\_\allowbreak{}lighton, music\_\allowbreak{}dislikeness, iot\_\allowbreak{}wemo\_\allowbreak{}on \\
MASSIVE-it & SetFit (mE5-base) & 0.787 & 0.000 & 0.000 & 9/59 & audio\_\allowbreak{}volume\_\allowbreak{}other, datetime\_\allowbreak{}convert, email\_\allowbreak{}addcontact \\
MASSIVE-it & mmBERT-base & 0.872 & 0.557 & 0.000 & 5/59 & iot\_\allowbreak{}hue\_\allowbreak{}lighton, music\_\allowbreak{}dislikeness, iot\_\allowbreak{}wemo\_\allowbreak{}on \\
MASSIVE-it & Qwen3-0.6B LoRA & 0.854 & 0.632 & 0.000 & 4/59 & iot\_\allowbreak{}hue\_\allowbreak{}lighton, music\_\allowbreak{}dislikeness, general\_\allowbreak{}greet \\
\midrule
Ottavo & Mammoth-1 & 1.000 & 0.978 & 0.947 & 0/24 & info\_\allowbreak{}prodotto\_\allowbreak{}prevendita, supporto\_\allowbreak{}uso\_\allowbreak{}prodotto, ritiro\_\allowbreak{}in\_\allowbreak{}negozio \\
Ottavo & Elephant-1 & 0.994 & 0.963 & 0.929 & 0/24 & supporto\_\allowbreak{}uso\_\allowbreak{}prodotto, info\_\allowbreak{}prodotto\_\allowbreak{}prevendita, programma\_\allowbreak{}fedelta \\
Ottavo & OpenAI Decisions & 1.000 & 0.964 & 0.947 & 0/24 & codice\_\allowbreak{}sconto, ritiro\_\allowbreak{}in\_\allowbreak{}negozio, fuori\_\allowbreak{}tema \\
Ottavo & TF-IDF + LR & 0.916 & 0.880 & 0.805 & 0/24 & fuori\_\allowbreak{}tema, info\_\allowbreak{}prodotto\_\allowbreak{}prevendita, ordine\_\allowbreak{}stato\_\allowbreak{}spedizione \\
Ottavo & SetFit (mE5-base) & 0.953 & 0.929 & 0.859 & 0/24 & fuori\_\allowbreak{}tema, rimborso\_\allowbreak{}stato, recesso\_\allowbreak{}ripensamento \\
Ottavo & mmBERT-base & 0.923 & 0.859 & 0.833 & 0/24 & recesso\_\allowbreak{}ripensamento, supporto\_\allowbreak{}uso\_\allowbreak{}prodotto, ritiro\_\allowbreak{}in\_\allowbreak{}negozio \\
Ottavo & Qwen3-0.6B LoRA & 0.918 & 0.854 & 0.824 & 0/24 & recesso\_\allowbreak{}ripensamento, articolo\_\allowbreak{}errato\_\allowbreak{}mancante, fuori\_\allowbreak{}tema \\
Ottavo & Qwen3-1.7B LoRA & 0.953 & 0.918 & 0.884 & 0/24 & fuori\_\allowbreak{}tema, recesso\_\allowbreak{}ripensamento, ordine\_\allowbreak{}stato\_\allowbreak{}spedizione \\
Ottavo & Qwen3-8B QLoRA & 0.980 & 0.951 & 0.941 & 0/24 & fuori\_\allowbreak{}tema, ritiro\_\allowbreak{}in\_\allowbreak{}negozio, supporto\_\allowbreak{}uso\_\allowbreak{}prodotto \\
\bottomrule
\end{tabular}
\end{table}
